\begin{textsample}{POS Dim 5 – global\_north\_2024\_07 – Score 22.00 – global\_north\_2024\_07/110987219.txt}  \label{ex:f5_pos_301}
NZ Must \textbf{Act} On ICJ Advisory Opinion On Israel On Friday, 19 July 2024, the \textbf{International} \textbf{Court} of Justice ( ICJ ) released a detailed and damning Advisory Opinion in relation to the ' \textbf{Legal} Consequences arising from the Policies and Practices of Israel in the Occupied Palestinian Territory, including East Jerusalem '.
There are three main sections below : i ) What the Advisory Opinion says, ii ) How New Zealand must respond ; and iii ) Links to further information. i ) What the Advisory Opinion says The Advisory Opinion is in response to a UN General Assembly ( UNGA ) Resolution adopted on 30 December 2022 which asked for the ICJ ' s \textbf{legal} opinion on two questions :"( a ) What are the \textbf{legal} consequences arising from the ongoing violation by Israel of the right of the Palestinian people to self-determination, from its prolonged occupation, settlement and annexation of the Palestinian territory occupied since 1967, including \textbf{measures} aimed at altering the demographic composition, character and status of the Holy City of Jerusalem, and from its adoption of related @ @ @ @ @ @ @ @ @ @ the policies and practices of Israel referred to... above affect the \textbf{legal} status of the occupation, and what are the \textbf{legal} consequences that arise for all States and the United Nations from this status?"( A/RES/77/247 ) The UNGA Request made it clear that Israel ' s conduct in relation to"the disastrous humanitarian situation and the critical socio-economic and security situation in the Gaza Strip"was to be considered by the \textbf{Court}, and the Advisory Opinion does lay out the reasons why it considers"effective control"over the Gaza Strip was maintained by Israel following its physical withdrawal in 2005.
However, the"apocalyptic"and"catastrophic"impact of Israel ' s conduct in Gaza since October 2023 is not covered by because that is subject to another ICJ \textbf{case} ( ' Application of the \textbf{Convention} on the \textbf{Prevention} and Punishment of the \textbf{Crime} of \textbf{Genocide} in the Gaza Strip ', see, for example, https : **39 ; 477 ; TOOLONG... ).
Are you getting our free newsletter?
In the Advisory Opinion, the @ @ @ @ @ @ @ @ @ @ * the State of Israel ' s continued presence in the Occupied Palestinian Territory is unlawful ; * the State of Israel is under an obligation to bring to an end its unlawful presence in the Occupied Palestinian Territory as rapidly as possible ; * the State of Israel is under an obligation to cease immediately all new settlement activities, and to evacuate all settlers from the Occupied Palestinian Territory ; and * the State of Israel has the obligation to make reparation for the damage caused to all the natural or \textbf{legal} persons concerned in the Occupied Palestinian Territory.
The \textbf{Court} also concluded that Israel is in breach of Article 3 of the \textbf{International} \textbf{Convention} on the Elimination of All Forms of \textbf{Racial} Discrimination ( ICERD ) -"States Parties particularly condemn \textbf{racial} segregation and apartheid and undertake to prevent, prohibit and eradicate all practices of this nature in territories under their \textbf{jurisdiction}"- although it did not specifically describe Israel as an apartheid state.In relation to the \textbf{legal} consequences for UN member states, the \textbf{Court} advised :"@ @ @ @ @ @ @ @ @ @ \textbf{legal} the situation arising from the unlawful presence of the State of Israel in the Occupied Palestinian Territory and not to render aid or assistance in maintaining the situation created by the continued. presence of the State of Israel in the Occupied Palestinian Territory".
The \textbf{Court} further advised :"\textbf{international} organizations, including the United Nations, are under an obligation not to recognize as \textbf{legal} the situation arising from the unlawful presence of the State of Israel in the Occupied Palestinian Territory ; and the United Nations, and especially the General Assembly, which requested the opinion, and the Security Council, should consider the precise modalities and further action required to bring to an end as rapidly as possible the unlawful presence of the State of Israel in the Occupied Palestinian Territory". ii ) How New Zealand must respond In the light of this Advisory Opinion confirming Israel ' s presence in the Occupied Palestinian Territories ( including Gaza ) is unlawful, the \textbf{Court} ' s advice that UN member states must not render aid or assistance in maintaining @ @ @ @ @ @ @ @ @ @ breach - among other things - of four \textbf{Provisional} \textbf{Measures} Orders of the ICJ in relation to the \textbf{Genocide} \textbf{Convention}, New Zealand must now ( finally ) take action in three key areas : * Withdraw immediately from RIMPAC"the world ' s largest"maritime live-firing war games which includes Israeli military personnel ( for further information see https : **39 ; 518 ; TOOLONG... ), \textbf{commit} to no military \textbf{cooperation} or participation in military exercises involving Israel, and convey New Zealand ' s deepest concern to the diplomatic representatives of the states that are supplying weapons and other military equipment to Israel ; * Impose \textbf{sanctions} on Israel in the same way as it has on Russia, which target individuals and entities that are of economic or strategic importance to Russia and include restrictions on trade in goods and services, restrictions on engaging in commercial activities, targeted financial \textbf{sanctions} and travel bans on individuals.
As a start, \textbf{sanctions} must be imposed on the political leaders who are giving the orders in relation to the Occupied Palestinian Territories ( including Gaza ) and @ @ @ @ @ @ @ @ @ @ carrying out those orders ; and * Formally \textbf{recognise} the State of Palestine without any further delay, as 145+ UN member states have already done.
Did you know \textbf{Scoop} has an Ethical Paywall?
If you ' re using \textbf{Scoop} for work, your organisation needs to pay a small \textbf{license} \textbf{fee} with \textbf{Scoop} \textbf{Pro}.
We think that ' s \textbf{fair}, because your organisation is benefiting from using our news resources.
In return, we ' ll also give your team access to \textbf{pro} news tools and keep \textbf{Scoop} free for personal use, because public access to news is important!

% matched lemmas: act, case, commit, convention, cooperation, court, crime, fair, fee, genocide, international, jurisdiction, legal, license, measure, prevention, pro, provisional, racial, recognise, sanction, scoop
\end{textsample}
